\chapter{Rozszerzenia systemu liczb rzeczywistych}
\label{ma1:ch:real-extensions}

Liczby rzeczywiste tworzą zupełne ciało uporządkowane. W analizie przydają się również obiekty opisujące nieograniczoność oraz wielkości nieskończenie małe. Przedstawimy dwie odmienne konstrukcje: rozszerzoną prostą rzeczywistą, która ma kresy wszystkich podzbiorów, lecz nie jest ciałem, oraz ciało liczb hiperrealnych, które jest uporządkowane, lecz nie jest archimedesowe ani zupełne porządkowo.

\section{Rozszerzona prosta rzeczywista}
\begin{definicja}
\label{ma1:def:extended-real-line}
Rozszerzona prosta rzeczywista to zbiór $\overline{\R}=\R\cup\{-\infty,+\infty\}$, gdzie $-\infty,+\infty$ są dwoma różnymi nowymi elementami, nienależącymi do $\R$. Rozszerzamy porządek liczb rzeczywistych przez warunek $-\infty<x<+\infty$ dla każdego $x\in\R$ oraz $-\infty<+\infty$.
\end{definicja}
Porządek jest liniowy: na $\R$ jest liniowy, a nowe elementy są odpowiednio najmniejszy i największy.

\begin{twierdzenie}
\label{ma1:thm:extended-complete}
Każdy podzbiór $E\subseteq\overline{\R}$ ma kres górny i dolny w $\overline{\R}$. W szczególności $\sup\varnothing=-\infty$, $\inf\varnothing=+\infty$.
\end{twierdzenie}
\begin{proof}
Dla zbioru pustego każdy element jest ograniczeniem górnym, a najmniejszym z nich jest $-\infty$. Niech $E\ne\varnothing$. Jeśli $+\infty\in E$ lub $E\cap\R$ jest nieograniczony z góry w $\R$, żaden element mniejszy od $+\infty$ nie jest ograniczeniem górnym. Zatem $\sup E=+\infty$.
W pozostałych przypadkach, jeśli $E\cap\R=\varnothing$, to $E=\{-\infty\}$ i $\sup E=-\infty$. Jeżeli $E\cap\R\ne\varnothing$, jest to zbiór ograniczony z góry w $\R$. Z zupełności istnieje $s=\sup_{\R}(E\cap\R)$. Jest on ograniczeniem górnym całego $E$, ponieważ jedynym możliwym dodatkowym elementem jest $-\infty$. Każde ograniczenie górne $E$ jest albo $+\infty$, albo rzeczywistym ograniczeniem górnym $E\cap\R$, więc jest nie mniejsze od $s$. Stąd $\sup E=s$.
Negację rozszerzamy przez $-(+\infty)=-\infty$, $-(-\infty)=+\infty$. Jest to inwolucja odwracająca porządek; stąd $\inf E=-\sup(-E)$, również dla zbioru pustego.
\end{proof}
\begin{przyklad}
Mamy $\sup\N=+\infty$, $\inf\{-n:n\in\N\}=-\infty$ i $\sup\{1-1/n:n\in\N_+\}=1$. Ostatnia równość wynika z własności Archimedesa: dla każdego $\varepsilon>0$ w $\R$ istnieje $n\ge1$ z $1/n<\varepsilon$.
\end{przyklad}

\section{Działania częściowe}
\label{ma1:sec:extended-operations}
\begin{definicja}
\label{ma1:def:extended-operations}
Na $\overline{\R}$ zachowujemy zwykłe dodawanie i mnożenie dwóch liczb rzeczywistych. Dla każdego $x\in\R$ definiujemy $x+(\pm\infty)=\pm\infty$, a także $(+\infty)+(+\infty)=+\infty$ i $(-\infty)+(-\infty)=-\infty$. Dodawanie jest przemienne; suma $(+\infty)+(-\infty)$ pozostaje nieokreślona.
Jeżeli przynajmniej jeden czynnik jest nieskończony, a drugi jest niezerowy, iloczyn ma wartość $+\infty$ dla jednakowych znaków i $-\infty$ dla przeciwnych znaków. Iloczyny $0\cdot(\pm\infty)$ pozostają nieokreślone. Odejmowanie oznacza dodawanie elementu przeciwnego.
\end{definicja}
W teorii miary przyjmuje się niekiedy dodatkową konwencję $0\cdot(+\infty)=0$ dla funkcji nieujemnych. Nie jest to uniwersalna reguła dla wyrażeń nieoznaczonych; w tym rozdziale jej nie przyjmujemy.

\begin{przyklad}
Dla $n\in\N_+$ pary $(n,-n)$, $(n,1-n)$ i $(2n,-n)$ mają pierwszy składnik dowolnie duży dodatni, a drugi dowolnie mały ujemny dla dostatecznie dużego $n$. Ich sumy wynoszą odpowiednio $0$, $1$ i $n$. Podobnie pary $(1/n,n)$, $(1/n^2,n)$ i $(1/n,n^2)$ mają pierwszy czynnik dodatni i dowolnie mały, drugi dowolnie duży, lecz iloczyny wynoszą $1$, $1/n$ i $n$. Stwierdzenia te wynikają z własności Archimedesa i nie wymagają jeszcze teorii granic.
\end{przyklad}
\begin{stwierdzenie}
\label{ma1:prop:extended-not-field}
Nie istnieje struktura ciała uporządkowanego na $\overline{\R}$ rozszerzająca porządek i działania na $\R$.
\end{stwierdzenie}
\begin{proof}
W dowolnym ciele uporządkowanym $x<x+1$ dla każdego elementu $x$, ponieważ $0<1$ i dodanie $x$ zachowuje nierówność. Dla $x=+\infty$ przeczyłoby to istnieniu elementu największego.
\end{proof}

\section{Filtry i ultrafiltry}
\label{ma1:sec:ultrafilters}
Pojęcie filtru pojawia się również w Topologii ogólnej, w rozdziale o zbieżności topologicznej. Przypomnimy potrzebną definicję i wyprowadzimy własności ultrafiltru używane do konstrukcji liczb hiperrealnych.
\begin{definicja}
\label{ma1:def:ultrafilter}
Niech $I$ będzie niepustym zbiorem. Filtrem właściwym na $I$ nazywamy rodzinę $\mathcal F\subseteq\mathcal P(I)$ taką, że $I\in\mathcal F$, $\varnothing\notin\mathcal F$, dla $A,B\in\mathcal F$ zachodzi $A\cap B\in\mathcal F$, a dla $A\in\mathcal F$ i $A\subseteq B\subseteq I$ zachodzi $B\in\mathcal F$. Ultrafiltrem nazywamy filtr właściwy maksymalny względem inkluzji.
\end{definicja}
\begin{lemat}
\label{ma1:lem:ultrafilter-alternative}
Filtr właściwy $\mathcal U$ na $I$ jest ultrafiltrem wtedy i tylko wtedy, gdy dla każdego $A\subseteq I$ dokładnie jeden ze zbiorów $A$ i $I\setminus A$ należy do $\mathcal U$.
\end{lemat}
\begin{proof}
Oba zbiory nie mogą należeć do filtru właściwego, gdyż ich przecięcie jest puste. Załóżmy maksymalność $\mathcal U$ i $A\notin\mathcal U$. Gdyby każdy $B\in\mathcal U$ spełniał $B\cap A\ne\varnothing$, rodzina
\[
 \mathcal G=\{C\subseteq I:\exists B\in\mathcal U,\ B\cap A\subseteq C\}
\]
byłaby filtrem właściwym: zawiera $I$, nie zawiera $\varnothing$, jest domknięta w górę, a świadkowie $B_1,B_2$ dają świadka $B_1\cap B_2$ dla przecięcia. Zawiera $\mathcal U$ i $A$, więc byłaby jego właściwym rozszerzeniem. Sprzeczność dowodzi istnienia $B\in\mathcal U$ z $B\cap A=\varnothing$. Wtedy $B\subseteq I\setminus A$, skąd $I\setminus A\in\mathcal U$.
Odwrotnie, przy założonej alternatywie żaden filtr właściwy nie może zawierać $\mathcal U$ i dodatkowo $A\notin\mathcal U$, gdyż zawierałby już $I\setminus A$.
\end{proof}
\begin{wniosek}
\label{ma1:cor:ultrafilter-partition}
Jeżeli $I$ jest rozłączną sumą skończenie wielu zbiorów $A_1,\ldots,A_k$, gdzie $k\ge1$, to dokładnie jeden należy do ultrafiltru $\mathcal U$ na $I$.
\end{wniosek}
\begin{proof}
Dwa różne składniki nie mogą należeć do $\mathcal U$. Gdyby żaden nie należał, wszystkie dopełnienia należałyby do $\mathcal U$, a ich przecięcie byłoby puste.
\end{proof}
\begin{przyklad}
Dla $i_0\in I$ rodzina $\{A\subseteq I:i_0\in A\}$ jest ultrafiltrem, zwanym głównym. Na $\N_+$ rodzina zbiorów o skończonym dopełnieniu jest filtrem właściwym, zwanym filtrem Frécheta. Nie jest ultrafiltrem: ani zbiór liczb parzystych, ani jego dopełnienie nie mają skończonego dopełnienia.
\end{przyklad}
\begin{twierdzenie}[Istnienie ultrafiltru niegłównego]
\label{ma1:thm:free-ultrafilter}
Przyjmując aksjomat wyboru, istnieje ultrafiltr $\mathcal U$ na $\N_+$ zawierający filtr Frécheta. Żaden zbiór skończony nie należy do takiego ultrafiltru.
\end{twierdzenie}
\begin{proof}
Porządkujemy przez inkluzję wszystkie filtry właściwe zawierające filtr Frécheta. Zbiór ten jest niepusty. Dla niepustego łańcucha jego suma mnogościowa jest filtrem właściwym: dwa zbiory należące do sumy należą wspólnie do jednego filtru z łańcucha, więc ich przecięcie należy do sumy. Pozostałe aksjomaty przechodzą na sumę bezpośrednio. Pusty łańcuch ma ograniczenie górne, mianowicie filtr Frécheta. Z lematu Zorna istnieje element maksymalny $\mathcal U$. Jest ultrafiltrem, ponieważ każde właściwe rozszerzenie nadal zawierałoby filtr Frécheta. Jeśli $A$ jest skończony, to $\N_+\setminus A\in\mathcal U$, więc $A\notin\mathcal U$.
\end{proof}
Ustalamy odtąd jeden taki ultrafiltr $\mathcal U$ na $I=\N_+$. Nie zakładamy istnienia wyróżnionego kanonicznego ultrafiltru niegłównego.
